\chapter{引言}

\section{研究背景与意义}

软件中潜藏的安全漏洞一旦被利用，后果往往远超开发者的预期。2017年WannaCry勒索软件在短短几天内瘫痪了全球数十万台主机，四年后Log4j漏洞又让几乎整个Java生态陷入紧急修补，两起事件的经济损失合计达数十亿美元\cite{chen2024}。这类案例反复说明一个现实：漏洞发现得越晚，修复代价越高。

模糊测试（Fuzzing）是目前应对这一问题最直接的手段之一。其原理并不复杂，就是向目标程序持续注入随机或经过变异的输入，看程序是否出现崩溃或异常行为\cite{li2024}。这种简单而高效的方法在工业界已有充分验证：Google的OSS-Fuzz平台运行至今，累计为超过1000个开源项目捕获并修复了上万个安全漏洞\cite{liu2023}。

但模糊测试的效果很大程度上取决于模糊驱动的质量。模糊驱动是一段适配程序，负责把模糊引擎产生的随机字节流转化为被测库能接受的API调用\cite{zhang2024}。写好一个驱动并不容易，开发者需要深入理解目标库的接口语义和调用约束，这往往要花数小时甚至数天\cite{lyu2024}。而且OSS-Fuzz上的驱动几乎全靠人工编写，最终平均覆盖率也只有约30\%\cite{liu2023}。所以本文认为，如何自动生成高质量的模糊驱动，是提升模糊测试规模和深度的关键瓶颈。

近两年大语言模型的进步让人看到了另一种可能性。LLM在海量代码语料上预训练后，对代码补全、生成和修复等任务表现出了相当强的泛化能力\cite{chen2021a}。值得注意的是，Google开源安全团队已经在实际项目中验证了这条路径，他们用LLM生成的模糊驱动不仅提升了覆盖率，还在OpenSSL中首次发现了此前未知的安全漏洞\cite{chang2024}。换句话说，LLM辅助模糊驱动生成已经不只是学术探索，它正在产生真实的安全收益。

不过，现有的LLM驱动生成方法仍有明显的短板。本文观察到四个共性问题：覆盖率反馈的粒度太粗（往往只到函数或API组合级别），单模型承担了从规划到生成到修复的全部职责导致注意力分散，轮次之间缺乏系统性的知识传递，以及生成策略对不同类型的库缺乏区分。这些问题共同制约了驱动质量的进一步提升。本文的思路是将多智能体协作架构与细粒度覆盖率反馈结合，让不同智能体各司其职，用行级覆盖率数据驱动精准迭代，同时在槽位隔离的知识空间中累积经验。这一方案旨在提升模糊驱动自动生成的效率和质量，从而降低人工编写驱动的成本、扩大模糊测试的适用范围。

\section{国内外研究现状}

模糊驱动的自动生成是提升模糊测试规模与深度的关键课题。围绕这一目标，研究者走出了两条主要路径：一条依赖程序分析来合成驱动，另一条借助大语言模型来生成驱动。下面分别梳理两条路径的代表性工作，分析各自的优势和局限，然后说明本文的切入点与创新点。

\subsection{基于程序分析的模糊驱动生成}

早期工作大多从已有代码中"学习"API的使用方式。Google的FUDGE系统\cite{babic2019}是这个方向的开创者，它分析目标库的使用者代码，提取API调用模式并据此合成驱动。在Google内部部署中，FUDGE为数千个库生成了上千个驱动，其中超过200个被集成到持续模糊测试中，帮助修复了150多个缺陷。Ispoglou等人的FuzzGen\cite{ispoglou2020}把这个思路往前推了一步，它分析多个消费者程序来构建抽象API依赖图（A$^2$DG），从而推断出更丰富的调用序列组合。

一个有意思的思路来自Zhang等人的APICRAFT\cite{zhang2021}。他们面对的是闭源SDK，因为拿不到源码，因此只能综合利用头文件、二进制和运行痕迹等信息来收集API依赖，再通过多目标遗传算法组合生成驱动。APICRAFT在macOS SDK上的实验很亮眼，自动生成的驱动比人工版本平均提升了64\%的覆盖率，8个月持续测试中发现了142个新漏洞，54个获得了CVE编号。三星研究院的Jeong等人则换了个角度，他们的UTOPIA工具\cite{jeong2023}直接从单元测试用例中提取API调用序列。这个思路很巧妙：单元测试本身就是开发者精心编写的、验证过API规范用法的代码，天然避免了从消费者程序中提取时混入自定义逻辑的问题。UTOPIA在55个项目中从8000余个单元测试生成了约5000个驱动，发现了123个漏洞。

更近期的工作开始引入形式化方法。Zhang等人在USENIX Security 2023提出RUBICK\cite{zhang2023}，用主动自动机学习算法把API使用模式建模为确定性有穷自动机，生成的驱动比FuzzGen多覆盖了50.42\%的代码边。Yu等人的Hopper\cite{yu2024}则从类型系统出发，自动推断参数间的依赖关系和取值约束，不需要消费者代码就能合成驱动。

总的来看，程序分析方法的优势在于能保证生成的驱动遵循真实的API调用模式。但它们的局限也很明显：必须有足够的先验样本（消费者代码、测试用例或使用实例），面对复杂语义约束和闭源库时仍然力不从心。

\subsection{基于大模型的模糊驱动生成}

大语言模型的出现为驱动生成带来了新范式。相比程序分析，LLM方法的通用性更强，不局限于特定语言或接口。按反馈机制的精细程度，现有方法大致可以分为三个层次。

\textbf{无迭代反馈的单次生成。}最基础的做法是让LLM一次性生成驱动，不做后续改进。Deng等人的TitanFuzz\cite{deng2023b}证明了LLM可以零样本地为深度学习库生成模糊输入，Xia等人的Fuzz4All\cite{xia2024}进一步把这个能力推广到任意编程语言。虽然这些工作针对的是测试输入而非驱动程序，但它们验证了LLM理解API语义的能力。在驱动生成领域，Google的OSS-Fuzz-Gen\cite{liu2023}采用单次生成模式。Zhang等人\cite{zhang2024}做了一个很有价值的大规模评估：在30个C项目的86个API上，用5种LLM和6种提示策略生成了73万多个候选驱动，跑了3.75 CPU年的对照实验。结论显示LLM生成驱动有潜力，但单次生成的质量和人工优化的驱动还有差距。问题的根源在于缺乏反馈，因为单次生成过程是开环的，无法根据执行效果做针对性改进。

\textbf{粗粒度反馈的迭代生成。}为了克服开环的局限，Lyu等人的PromptFuzz\cite{lyu2024}引入了覆盖率反馈来指导迭代。它通过指令化程序生成、错误自动校验和覆盖率驱动的提示变异来探索未覆盖区域。在14个库上的测试中，PromptFuzz比OSS-Fuzz默认驱动和Hopper分别提高了1.61倍和1.63倍的分支覆盖率，还发现了33个新漏洞。但本文注意到一个关键限制，PromptFuzz的反馈停留在API组合级别：它知道哪些API组合被覆盖了，却不知道具体哪行代码没跑到，更无法诊断未覆盖的原因。这种粗粒度反馈让LLM只能做相对盲目的变异，难以高效突破覆盖率瓶颈。

\textbf{知识增强的生成方法。}Xu等人的CKGFuzzer\cite{xu2024}尝试用知识图谱来增强LLM对被测库的理解，在8个库上平均提升了8.73\%的覆盖率。但CKGFuzzer的知识图谱是静态构建的，不随迭代更新；其多智能体之间的协作也只是顺序串联，各智能体共享相同上下文，缺乏基于覆盖率反馈的差异化信息传递。

\subsection{现有方法的不足与本文切入点}

综合来看，程序分析方法可靠但依赖先验代码，LLM方法灵活但可控性不足。特别是在LLM方法中，本文认为当前研究存在四个共性短板：

第一，覆盖率反馈粒度不足。现有迭代方法的反馈多停留在API组合或函数级别，无法告诉LLM具体哪行代码没覆盖、为什么没覆盖。理想的反馈应该精确到行和分支，并附带根因分析。

第二，单一模型职责过载。让一个LLM同时做规划、生成、修复和分析，上下文一长注意力就分散了。Liu等人\cite{liu2024lost}的研究指出LLM处理长上下文时存在"中间遗忘"现象，即关注首尾而忽略中间。在代码生成场景中，把API签名、覆盖率数据、编译错误、历史代码全塞进一个提示里，模型对关键信息的利用效率会明显下降。这个问题不是简单拆分提示词就能解决的，因为各环节需要不同的知识视图、独立的温度控制和隔离的状态管理，这些需求天然指向多智能体架构。

第三，缺乏结构化知识积累。多数方法每轮迭代相互独立，成功的代码模式、失败的尝试方向、发现的API约束都没有被系统性地记录和复用。后续轮次可能重蹈覆辙，也可能遗忘已验证有效的模式。

第四，生成策略单一。用统一的提示词模板生成所有驱动，忽略了不同类型库的差异。解析类库需要把原始字节直接喂给解析函数，有状态库则需要构造多步调用序列，两种场景的代码结构、输入分配和资源管理完全不同，单一提示词无法兼顾。

本文正是针对这四个不足展开工作：用多智能体架构解决注意力稀释，用行级覆盖率反馈实现精准迭代引导，用按槽位隔离的知识累积实现经验传递，用策略库加自动匹配实现对不同库的自适应生成。

\section{研究目标}

本文希望解决的核心问题是：如何让LLM生成的模糊驱动在反馈精度、职责分配、知识积累和策略适配这四个维度上都有实质性的提升？为此，本文设计并实现了MAGIC4FDG（Multi-Agent Guided Iteration with Coverage for Fuzz Driver Generation），一个端到端的模糊驱动自动生成系统。具体目标从四个维度展开。

在架构层面，本文需要实现由七个专职智能体组成的有向无环图协作架构。每个智能体只负责一件事：知识提取、策略规划、代码生成、编译修复、覆盖率分析或缺口诊断，并且只接收与自身职责相关的知识子集。这样做的目的是避免单模型上下文过载导致的注意力稀释。

在反馈层面，本文需要构建基于llvm-cov的行级与分支级覆盖率反馈闭环。不只是告诉系统"覆盖率是多少"，而是通过缺口诊断智能体分析具体哪些代码行没覆盖、为什么没覆盖，把分析结果转化为结构化约束，给代码生成智能体一个明确的改进方向。

在迭代层面，本文需要设计探索-利用双阶段范式。第一轮通过LLM场景推理和多样化策略库做广度探索，建立多条独立的生成轨迹；后续轮次基于约束反馈对各轨迹独立做增量改进。每条轨迹维护隔离的知识空间，跨轮次累积约束和有效模式，形成持续增长的经验库。

在评测层面，本文需要在OSS-Fuzz收录的多个开源C/C++库上做系统性实验，与PromptFuzz、OSS-Fuzz-Gen等工具和人工驱动全面对比，并通过消融实验验证各组件的贡献。

\section{研究内容}

\subsection{主要工作}

围绕上述目标，本文工作聚焦于MAGIC4FDG系统的体系化设计与工程实现，核心要回答三个问题：多智能体协作如何提升生成质量？覆盖率反馈如何实现精准迭代？知识累积如何提升迭代效率？具体包括四个方面。

第一，多智能体协作架构的设计与实现。本文将驱动生成拆成七个子任务，每个交给一个专职智能体，用DAG调度框架按拓扑顺序执行。其中关键设计是差异化上下文注入，即同一个知识库，不同智能体获取的上下文不一样。举例来说：策略规划智能体看到的是宏观的API分类摘要和覆盖率趋势，代码生成智能体看到的是具体的目标API签名和累积约束，缺口诊断智能体看到的是完整调用图和未覆盖的源码行。这样做的原因是如果把所有信息一并提供给每个智能体，注意力就会分散，反而降低生成质量。另外系统还设计了条件路由：变体全部编译失败就跳过模糊测试节省资源，覆盖率连续几轮不涨就自动终止。

第二，多提示词策略自适应生成范式。本文搭建了一个包含十种模糊策略的策略库，包括解析驱动、多API序列、往返验证、结构感知等，每种策略有YAML元数据描述适用条件，还有策略正文作为提示词后缀注入LLM。同一轮里不同槽位拿到的基础模板一样，但策略后缀不同，这样产出驱动的探索风格就会不一样。探索阶段Planner智能体通过场景推理做策略匹配，如果推理失败还有基于特征标签的规则引擎兜底。在知识提取模块，本文用clang AST实现了五步流程：语法树解析、宏解析、调用图构建、语义标注和分类推断。

第三，覆盖率反馈迭代与知识累积机制。覆盖率数据的采集依赖llvm-cov，精确到源码行级别。模糊测试本身在Docker容器内以fork模式运行（这样单个输入触发的崩溃不会中断整个测试进程），测试结束后再通过语料库重放来收集完整的profraw数据。每条生成轨迹拥有独立的知识空间，随轮次推进不断累积正负模式，比如"必须先调用cJSON\_Parse才能操作返回的对象"属于正模式，"传入零长度缓冲区会触发断言失败"则是负模式。缺口诊断智能体的职责被刻意限定为只做分析、不写代码，它需要指出哪些行没覆盖并给出可能的原因，输出结构化约束交给生成智能体去落实。这种诊断与生成的分离是有意为之的，混在一起容易导致角色混淆、输出质量下降。收敛判定分两个层级：某条轨迹连续若干轮覆盖率不涨，就尝试更换策略；如果所有轨迹的整体涨幅都停滞了，则终止任务避免浪费算力。

第四，系统工程化实现与实验评估。工程层面，所有执行环节（编译、模糊测试、覆盖率采集）都在容器内完成，宿主机只负责调度。构建缓存机制让同一个库只需编译一次，后续轮次直接挂载缓存产物。效率方面做了多层并行：LLM调用通过异步API轮询实现并发，多个变体的编译修复和模糊测试也是并行执行的。实验在12个OSS-Fuzz收录的C/C++库上展开，对比对象包括PromptFuzz、OSS-Fuzz-Gen以及各库已有的人工驱动，同时设计了消融实验来量化各组件（多智能体拆分、行级反馈、知识累积、策略库）的独立贡献。

\subsection{论文结构}

全文共六章，组织如下：

第一章是引言，交代研究背景和动机，梳理国内外研究现状，指出现有方法的不足，明确研究目标和内容。

第二章介绍相关技术基础，包括模糊测试与LibFuzzer引擎、覆盖率引导技术、大语言模型代码生成、多智能体协作框架以及clang AST静态分析，为后续设计提供技术铺垫。

第三章是MAGIC4FDG的需求分析与概要设计。先用系统流程图展示整体运行逻辑，再做涉众分析、功能性与非功能性需求分析和用例建模，最后通过"4+1"架构视图完成概要设计，明确模块职责和协作关系。

第四章是详细设计与实现。按模块逐一阐述各智能体的设计方案和核心实现，每个模块包含类图、顺序图和关键代码片段。

第五章是实验设计与结果分析。说明实验设置后，围绕系统性能和组件有效性进行多组实验，并通过长时模糊测试验证生成驱动的缺陷发现能力，用定量数据验证系统。

第六章总结全文，分析局限性，并展望未来该领域可能的方向。